\documentclass[11pt]{article}
\usepackage[margin=0.85in]{geometry}
\usepackage[T1]{fontenc}
\usepackage{amsmath,amssymb,booktabs,longtable,microtype,xurl,hyperref}
\setlength{\emergencystretch}{3em}
\hypersetup{colorlinks=true,urlcolor=blue,linkcolor=blue}
\title{Worst-Frame Runtime: Diagnosis and a Measured Route Toward Fast PCBF/CLF Control}
\author{Pan Dynamics Collision Avoidance Research}
\date{October 2, 2026}
\begin{document}\maketitle

\section{Finding and scope}
The longest recorded call is the 15 m/s braking-lead encounter at 5.90 s:
2.952268 s. Its first completed PCBF/CLF pair already meets the implemented
prediction-accuracy thresholds and predicts a collision-free frozen-input
continuation in a subsequent dense numerical audit. Seven additional pairs
are not needed to obtain that first observed collision-free candidate. They
are triggered by positive CLF slack at the numerical first-input correction
boundary. Expensive repeated convex solves, rather than flow construction
or nonlinear prediction evaluation, dominate the call.

The investigation exactly replays all 119 calls through that hold, with zero
maximum applied-input difference from the original campaign. It captures the
16 convex subproblems, profiles the complete call, and benchmarks isolated
solver and iteration changes. Production controller source is unchanged.
These are diagnostic experiments and design recommendations, not an
announcement that the online controller now meets 0.1 s or that all scenarios
are safe. The original crossing intersample collision remains unresolved.

\section{Measured cost and structure}
The original unprofiled call has these disjoint components:
\begin{center}\begin{tabular}{lr}\toprule Component & Seconds\\\midrule
PCBF solves & 0.589896\\
CLF solves & 2.149130\\
Formulation excluding CLF construction & 0.096662\\
CLF construction & 0.092109\\
Nonlinear prediction agreement & 0.015625\\
Initialization & 0.000699\\
Other bookkeeping & 0.008147\\\bottomrule\end{tabular}\end{center}
CLF solves account for approximately 72.8\% of the call; the two solver
stages together account for 92.8\%. Removing prediction evaluation would
save only about 0.5\%, while removing the mechanism that controls model error.
The source already uses compiled bicycle and nominal-value kernels.

The horizon contains 76 holds: 16 prefix holds and 60 completion holds,
covering 3.8 s. Each primary has 630 decision variables, 462 dynamic
initial-condition/equality rows, and 1069 linear inequality rows before
variable bounds. The secondary has 632 variables, 1070 such inequalities,
and second-order cones of dimensions 6, 6 and 154. The large cone represents
the squared input-increment regularizer. Both phases start fresh numerical
solves; shifting a nonlinear reference trajectory is not reuse of an
interior-point iterate or factorization.

The exact replay reports 137 primary and 242 secondary interior-point
iterations, 379 in total. The observed linear-system method is
\texttt{augmented} in MATLAB R2026a Update 3. Profiling attributes 2.867 s
of its instrumented call to \texttt{coneprog}/\texttt{interiorPointMethod};
profile overhead means this number must not replace the original runtime.
The nominal CLF kernel is called 88 times and takes approximately 0.101 s
under the profiler. It is a secondary optimization target, not the main cause
of the multi-second call.

The CLF successor value depends on only the first two inputs, but those inputs
must admit a constrained future trajectory. It is consequently not valid to
discard the completion variables and solve an unconstrained two-variable QP.
The endpoint ellipsoid and CLF epigraph are quadratic constraints; the whole
second stage is not a linearly constrained QP merely because its local value
model is quadratic.

\section{Why the eight rounds are repeated}
The first round's maximum body-pose prediction discrepancy is 1.693 mm,
below the configured 10 mm threshold. Every later round also meets the
threshold. The inherited input trust scale is 0.138209, so its steering
correction is limited to about $\pm0.010366$ rad around the current anchor.
The same small radius is retained within the frame even after agreement is
excellent; its growth rule applies only after a frame has been accepted.

\begin{center}\small\begin{tabular}{rrrrr}\toprule
Round & PCBF optimum & Modeled CLF slack & Body error (mm) & Solve time (s)\\\midrule
1 & $1.02\,10^{-8}$ & 68.8788 & 1.693 & 0.3538\\
2 & $8.08\,10^{-9}$ & 68.2901 & 0.918 & 0.3282\\
3 & $4.62\,10^{-9}$ & 67.9730 & 0.456 & 0.3135\\
4 & $7.80\,10^{-10}$ & 67.7911 & 0.074 & 0.3374\\
5 & $7.07\,10^{-10}$ & 67.6795 & 0.046 & 0.3421\\
6 & $8.76\,10^{-9}$ & 67.6058 & 0.158 & 0.3420\\
7 & $4.72\,10^{-9}$ & 67.5541 & 0.278 & 0.3174\\
8 & $1.50\,10^{-7}$ & 67.5540 & 0.011 & 0.4047\\\bottomrule\end{tabular}\end{center}
The retained candidate is round seven. The total modeled CLF-slack
improvement over round one is about 1.325, or 1.92\%. Round eight no longer
hits the first-input correction boundary and does not improve the retained
nonlinear result. Positive slack alone does not establish that a further
iteration can achieve zero slack or is necessary for safety.

An independent ODE45 rollout of each complete frozen input sequence uses
relative tolerance $10^{-11}$, absolute tolerance $10^{-12}$, and 31 output
samples per 50 ms hold. Round one's minimum observed clearance over all
76 holds is 0.099685 m; round seven's is 0.099241 m. Their first-hold minima
are 0.099843 and 0.099771 m. These observed clearances do not constitute a
continuous-time certificate or prove recovery after switching the entire
closed-loop implementation to one round.

\section{Exact primary optimality without a numerical solve}
Let the primary objective be $J_s=\sum_i\xi_i$, with $\xi_i\ge0$. If the current
anchor satisfies every constraint with $d=0$ and $\xi=0$, then
\[
 0\le J_s^\star\le J_s(0,0)=0.
\]
The primary optimum is known exactly in the mathematical model. A new
interior-point solve provides no additional optimality information. The
secondary CLF problem still has to be solved.

In captured rounds two through seven, the zero correction satisfies all
linear inequalities, equality defects, bounds and the endpoint cone. Thus
six of the eight primary optimizations admit this certificate. Round one's
anchor violates the terminal cone by approximately $9.81\,10^{-4}$, so it
cannot use this shortcut. Round eight has a positive linear residual of
approximately $1.26\,10^{-6}$ and also lacks this zero-anchor certificate.
The certificate concerns the complete convex primary, not just the collision
rows, and is not a nonlinear safety certificate.

An isolated implementation of this exact test reduces numerical conic calls
from 16 to 10. It continues to solve the CLF stage in every round. This is
algebraic completion of the primary optimization, not a backup policy or
omission of the second objective.

\section{Equivalent convex expressions and solver measurements}
The current secondary objective is
\[
 \min\ \bar\rho+\varepsilon\sigma,\qquad
 \|W\Delta U\|_2^2\le\sigma\le1.
\]
Its input boxes imply $\|W\Delta U\|_2^2\le1$, with the normalization already
used by the controller. Therefore the epigraph variable and its 154-dimensional
cone can be eliminated exactly in the projected feasible set:
\[
 \min\ \bar\rho+\varepsilon\Delta U^T W^T W\Delta U.
\]
The remaining endpoint and CLF cones stay. A solver supporting quadratic
objectives can represent this expression directly; the PCBF priority and
penalty relative to the linearization inputs are unchanged.

All 16 captured problems were solved five times per representation through
the compiled Clarabel 0.11.1 Python binding, using one thread. Each repetition
creates a fresh solver and includes matrix conversion, construction and solve
time. No warm-start benefit is assumed.
\begin{center}\small\begin{tabular}{lrr}\toprule
Representation & Sum of per-problem medians (s) & Solve only (s)\\\midrule
Same conic epigraphs & 0.367088 & 0.307370\\
Native quadratic regularizer & 0.307435 & 0.251810\\\bottomrule\end{tabular}\end{center}
Every solve in these two comparisons returns \texttt{Solved}. The largest
residual after reconstructing the original constraints is
$1.23\,10^{-11}$. The largest absolute objective difference from the captured
MATLAB points is $1.18\,10^{-6}$ in their original scaled convex objectives.
These are numerical agreement results, not bitwise-identical optimizer
trajectories.

The original solver sum was 2.739026 s. The 0.307435 s alternative is a
sum of medians for fixed captured problems, not an end-to-end new controller
measurement. Even adding the original non-solver time optimistically would
leave roughly 0.52 s. Solver replacement alone cannot establish a 0.1 s
controller deadline.

Other measured alternatives were unfavorable. Selecting MATLAB
\texttt{normal} or \texttt{schur}, or converting the small dense cone data
to sparse storage, did not provide a comparable speedup; some final
secondary solves returned a stalled status. Full state elimination reduced
variables but densified long-horizon constraints and increased representative
solves to approximately 1.03--1.14 s. Blindly reducing the number of variables
is not the same as exploiting the dynamic block structure.

A further diagnostic fixed all prefix slacks to zero. This is justified only
when zero primary optimality has been established. The final captured model
instead returns \texttt{PrimalInfeasible}; its finite numerical certificate
must not be treated as a control trajectory. Thus the observed near-zero
primary values do not justify universally dropping all safety slack variables.
For the certified zero cases, eliminating them expresses the exact
lexicographic optimum, but differs from the existing $10^{-6}$ tie allowance.

\section{Measured full-frame counterfactuals}
Each cached prototype was evaluated three times at the identical measured
state and previous-plan fixture. These timings include the controller call;
offline audits are excluded.
\begin{center}\small\begin{tabular}{lrrrr}\toprule
Prototype & Median (s) & Rounds & Conic calls & CLF slack\\\midrule
Original & 3.052923 & 8 & 16 & 67.5541\\
Certified primary shortcut & 2.548748 & 8 & 10 & 67.5542\\
Shortcut and adaptive radius & 1.658542 & 5 & 8 & 67.4694\\
Stop at first accurate complete pair & 0.408166 & 1 & 2 & 68.8788\\\bottomrule\end{tabular}\end{center}
The adaptive-radius diagnostic expands the input correction scale within the
frame when the prediction-error ratio is at most 0.25 and the first-input
box is active, using $\min(2,0.8/\sqrt r)$ and a maximum nominal scale of one.
It retains ordinary accuracy checks and contracts inaccurate steps as before.
It produces slightly lower CLF slack in five rounds; its frozen-input minimum
observed clearance is 0.098142 m. This is evidence of an avoidable small-step
obstruction in this frame, not a uniform iteration bound.

The one-round diagnostic is deliberately not adopted. Earlier closed-loop
experiments showed that stopping at an accurate but bound-limited CLF step
can repeatedly shift to an unhelpful future-input reference and prevent
nominal recovery. A fast safe first hold does not establish successful
long-term recovery. Each local convex CLF problem must still be optimized;
what needs redesign is when to rebuild and solve the entire problem again.

\section{Recommended mathematical design}
The evidence supports one lexicographic controller with the following changes,
in this order:
\begin{enumerate}
\item Preserve both objectives and exploit known lower bounds. Use a feasible
zero-slack primal point to complete the PCBF minimization analytically. If that
certificate is unavailable, solve the primary normally. Never infer optimality
or safety from a finite failed-solver array.
\item Use the native quadratic regularizer and a compiled sparse conic solver.
Preserve the dynamic block structure. Reuse fixed-size matrix storage and
symbolic structure when possible; do not conflate this with an unsupported
primal/dual warm-start interface. A C/MEX integration needs its own complete
controller benchmarks.
\item Adapt the trust radius during the frame using actual versus predicted
model agreement. Positive CLF slack and contact with an artificial box are not
sufficient reasons to repeat up to eight small identical steps. Examine
constraint activity and useful predicted improvement. After an accurate
completed pair, spend further work only on repairing a material model defect
or obtaining meaningful additional dissipation. A stopping rule must be
validated specifically against the earlier loss-of-recovery regressions.
\item Couple optimization accuracy to control accuracy. Keep feasibility
accuracy distinct from objective optimality. Near nominal behavior, the
allowed CLF/model error must shrink relative to the required Lyapunov decrease;
a fixed loose slack tolerance would only support practical convergence.
The current CLF model permits a 0.01 scaled prediction discrepancy while its
convex solver targets an approximately $10^{-8}$ optimality tolerance. This
mismatch motivates a derived inexact-solve criterion, not arbitrary loosening
of collision feasibility.
\item Close the temporal safety gap with interval-aware geometry. Faster
node optimization cannot repair the separately observed crossing collision.
For an available bound $L$ on distance variation over a constraint interval,
\[
 d(t)\ge\min\{d(t_j),d(t_{j+1})\}-L\Delta t/2.
\]
Model error and safety slack must also be covered by the available margin.
Conservative bounds or local temporal refinement are needed where this reserve
is insufficient. Critical-interval rows can be refined without uniformly
multiplying the entire horizon's state and input variables. The bound must be
valid over the trial region; sampled velocities alone are not such a proof.
\end{enumerate}

For the first captured pair, the native-quadratic convex calls cost about
0.038 s including setup, compared with approximately 0.354 s of original
solver time. Combined with roughly 0.025 s of formulation and other work,
this suggests that a useful complete pair within 0.1 s is plausible in this
frame. It is an estimate from separately measured components, not a measured
integrated latency, and it does not cover repeated restoration or cold starts.
A uniformly successful two-round policy has not been demonstrated.

The next implementation experiment should first integrate the two equivalent
changes and measure all fourteen closed-loop cases with the existing nonlinear
accuracy rules. Then evaluate the stopping/trust update jointly with recovery
and interval clearance. Acceptance requires collision avoidance, eventual
nominal recovery, and measured full-call latency, including difficult frames
and initialization. The 50 ms nominal sampling period and any actual computation
delay must also be represented consistently. No backup controller or second
CLF is required by the proposed changes. A fixed iteration cap alone proves
neither existence of a safe returned solution nor a wall-clock deadline.

\section{Reproducibility and references}
The baseline is project commit
\path{660a4cdd0875c8650665609ff48bb8cc746e1409}. The adjacent evidence directory
contains round diagnostics, all benchmark repetitions, offline audits, and
source/input hashes. The exact fixture, captured matrices, profiler data,
helper source and cached experimental solver copies are retained in the
external experiment export. The isolated Python environment and solver
vendor tree are excluded. No production implementation change or new full
closed-loop campaign is included in this analysis. No MATLAB unit-test rerun
was needed for an unchanged controller.

Complete recorded validation includes 119 exact replay calls, 108 MATLAB and
240 Clarabel frozen-problem benchmark calls, twelve timed prototype controller
calls, and frozen-input ODE45 audits. Solver failures and the infeasible
zero-slack experiment are retained rather than counted as successful controls.
All measurements use a desktop AMD Ryzen 7 7800X3D host; MATLAB uses
\texttt{-singleCompThread}. The host is not an isolated real-time target.

The official \href{https://www.mathworks.com/help/optim/ug/coneprog.html}{MATLAB
coneprog documentation} describes the conic interface and linear-system
choices. Its live page currently describes R2026b; this study's actual
solver choices and outputs come from the installed R2026a implementation.
The official \href{https://clarabel.org/}{Clarabel documentation} describes
quadratic-objective conic problems. Its
\href{https://clarabel.org/stable/user_guide_data_updating/}{data-update guide}
requires fixed dimensions and sparsity, unchanged cone definitions, and
compatible preprocessing settings. Such data updates are not evidence that
an arbitrary previous primal/dual iterate is accepted as a warm start.
\end{document}
